% ============================================================
% 05_results.tex  (revised 2026-10; numbers from results/power_summary and the run tables)
% ============================================================

\subsection{Point forecasts}
\label{subsec:res_point}

Tables~\ref{tab:yulara} and~\ref{tab:nist} report the test errors in power. At Yulara the values
are means and standard deviations over five seeds; at NIST a single seed is reported.

\begin{table}[t]
  \centering
  \caption{Yulara (four uncurtailed sites, 768~kW), test period 2019-07-19 -- 2020-05-13, lead
  times 1--24~h, daytime targets. Mean $\pm$ standard deviation over five seeds. Skill: MSE
  skill score against day-ahead persistence.}
  \label{tab:yulara}
  \small
  \begin{tabular}{lcccccc}
    \toprule
    & \multicolumn{3}{c}{without GFS} & \multicolumn{3}{c}{with GFS} \\
    \cmidrule(lr){2-4}\cmidrule(lr){5-7}
    Model & MAE (kW) & nRMSE (\%) & Skill & MAE (kW) & nRMSE (\%) & Skill \\
    \midrule
    AHSE                   & \textbf{36.18 $\pm$ 0.32} & \textbf{8.21} & \textbf{0.36} & \textbf{31.32 $\pm$ 0.05} & \textbf{6.95} & \textbf{0.54} \\
    AHSE-global            & 37.37 $\pm$ 0.48 & 8.23 & 0.36 & 31.34 $\pm$ 0.06 & 6.95 & 0.54 \\
    XGBoost                & 41.35 $\pm$ 0.12 & 8.61 & 0.30 & 31.34 $\pm$ 0.06 & 6.95 & 0.54 \\
    LightGBM               & 42.07 $\pm$ 0.07 & 8.76 & 0.27 & 32.36 $\pm$ 0.03 & 7.08 & 0.53 \\
    PatchTST               & 40.02 $\pm$ 0.89 & 8.77 & 0.27 & -- & -- & -- \\
    N-HiTS                 & 40.62 $\pm$ 1.17 & 8.86 & 0.25 & -- & -- & -- \\
    GRU                    & 40.41 $\pm$ 0.94 & 8.62 & 0.30 & -- & -- & -- \\
    LSTM                   & 42.50 $\pm$ 1.77 & 9.10 & 0.22 & -- & -- & -- \\
    GFS clearness index    & -- & -- & -- & 46.54 & 9.58 & 0.13 \\
    Day-ahead persistence  & 40.56 & 10.27 & 0 & 40.56 & 10.27 & 0 \\
    \bottomrule
  \end{tabular}

  \smallskip
  \footnotesize The neural members do not use GFS; their errors are identical in both columns.
\end{table}

\begin{table}[t]
  \centering
  \caption{NIST Ground array (271~kW), test period 2017-06-04 -- 2017-10-10, lead times
  1--24~h, daytime targets, seed 42.}
  \label{tab:nist}
  \small
  \begin{tabular}{lcccccc}
    \toprule
    & \multicolumn{3}{c}{without GFS} & \multicolumn{3}{c}{with GFS} \\
    \cmidrule(lr){2-4}\cmidrule(lr){5-7}
    Model & MAE (kW) & nRMSE (\%) & Skill & MAE (kW) & nRMSE (\%) & Skill \\
    \midrule
    AHSE                   & \textbf{27.55} & \textbf{14.52} & \textbf{0.41} & \textbf{20.39} & \textbf{10.80} & \textbf{0.68} \\
    AHSE-global, XGBoost   & 27.84 & 14.79 & 0.39 & 20.47 & 10.81 & 0.68 \\
    LightGBM               & 28.47 & 15.26 & 0.35 & 20.79 & 11.06 & 0.66 \\
    GRU                    & 28.89 & 14.75 & 0.39 & -- & -- & -- \\
    LSTM                   & 29.67 & 15.06 & 0.37 & -- & -- & -- \\
    PatchTST               & 29.97 & 15.87 & 0.30 & -- & -- & -- \\
    N-HiTS                 & 30.51 & 15.67 & 0.32 & -- & -- & -- \\
    GFS clearness index    & -- & -- & -- & 28.63 & 14.26 & 0.43 \\
    Day-ahead persistence  & 34.13 & 18.96 & 0 & 34.13 & 18.96 & 0 \\
    \bottomrule
  \end{tabular}
\end{table}

\paragraph{Weather forecasts are the main source of skill.}
Adding GFS covariates lowers the MAE of AHSE from 36.18 to 31.32~kW at Yulara (nMAE 4.71\%
$\to$ 4.08\%) and from 27.55 to 20.39~kW at NIST (10.17\% $\to$ 7.52\%). The gain is significant
for every seed at Yulara (4.4 to 5.2~kW; Diebold--Mariano statistic 3.2--3.8 on daily losses) and
at NIST (7.17~kW, 95\% interval [5.51, 8.92]~kW, DM 5.7). It is larger in the humid subtropical
climate of Maryland, where the clearness index is far less persistent from one day to the next
than in the desert, and it grows with the lead time (Fig.~\ref{fig:leadtime}): past observations
only inform the first hours, whereas the weather forecast carries information for the whole
horizon. Used alone, the GFS clearness index is worse than day-ahead persistence at Yulara and
better at NIST; its value comes from the learned combination with recent observations.

\begin{figure}[t]
  \centering
  \includegraphics[width=\linewidth]{fig_mae_leadtime.pdf}
  \caption{Test MAE in kW by lead time, without GFS (dashed) and with GFS (solid), seed 42. Day-ahead
  persistence is shown for reference.}
  \label{fig:leadtime}
\end{figure}

\paragraph{AHSE combines models only when the evidence supports it.}
Without GFS at Yulara, the selector keeps N-HiTS as reference and combines models on every block
(LightGBM, XGBoost and N-HiTS on the first hour; four models on 1--6~h; day-ahead persistence,
N-HiTS and XGBoost beyond 6~h). AHSE then beats the best single model of every seed, which is
N-HiTS, GRU or PatchTST depending on the seed, by 2.8 to 3.9~kW (7 to 10\%; DM 5.0--7.2), and
AHSE-global by 1.2~kW on average. With GFS, the GFS-fed XGBoost dominates the pool: the guard keeps
it alone from the second hour on, and AHSE equals it (differences of at most 0.05~kW). At NIST,
AHSE is not significantly different from XGBoost, with or without GFS (gains of 0.28 and
0.08~kW, $p = 0.22$ and $0.48$), and is significantly better than LightGBM. AHSE is therefore
never significantly worse than the best member chosen in hindsight on the test set; it improves
on it when the pool contains complementary members and falls back to it otherwise. The seed
variability of AHSE (0.32~kW without GFS, 0.05~kW with GFS) is an order of magnitude smaller
than these differences.

\paragraph{A metric caveat.}
At Yulara without GFS, day-ahead persistence has a lower MAE (40.6~kW) than XGBoost and LightGBM,
but a much larger RMSE (10.3\% against 8.6\% of capacity): in a sunny desert the clearness index
is very persistent and the MAE rewards forecasts that are exact on most days and badly wrong on
the few cloudy ones. Conversely, a MAE computed on $k_t$ instead of power gives the same weight
to dawn and dusk hours as to noon; on that scale AHSE appeared significantly behind XGBoost at
NIST without GFS, a difference that disappears in power. Conclusions should therefore rest on
power errors and include a squared-error measure.

% ------------------------------------------------------------
\subsection{Prediction intervals}
\label{subsec:res_intervals}

\begin{table}[t]
  \centering
  \caption{90\% prediction intervals of AHSE on the test period ($k_t$ units; out-of-fold
  calibration on the validation period). IS: interval score (lower is better).}
  \label{tab:intervals}
  \small
  \begin{tabular}{llccc}
    \toprule
    Configuration & Method & Coverage & Width & IS \\
    \midrule
    Yulara, no GFS & split, pooled            & 0.909 & 0.604 & 0.864 \\
                   & split, Mondrian          & 0.910 & 0.525 & 0.800 \\
                   & adaptive (ACI), Mondrian & 0.901 & 0.495 & 0.805 \\
    Yulara, GFS    & split, pooled            & 0.903 & 0.491 & 0.728 \\
                   & split, Mondrian          & 0.907 & 0.420 & 0.657 \\
                   & adaptive (ACI), Mondrian & 0.902 & 0.413 & 0.656 \\
                   & split, binned by GFS     & 0.916 & 0.440 & \textbf{0.627} \\
    NIST, GFS      & split, Mondrian          & 0.921 & 0.790 & 1.002 \\
                   & adaptive (ACI), Mondrian & 0.910 & 0.750 & 0.998 \\
                   & split, binned by GFS     & 0.907 & 0.690 & \textbf{0.949} \\
    \bottomrule
  \end{tabular}
\end{table}

All calibrated intervals reach their nominal coverage on the test period
(Table~\ref{tab:intervals}). At Yulara, Mondrian calibration by lead-time block and forecast
level lowers the interval score by 7--10\% against pooled calibration and adaptive conformal
inference brings the coverage closest to the nominal level; at NIST the Mondrian cells do not
improve on pooled calibration (interval score 1.002 against 0.998). GFS narrows the 90\%
intervals by 17\% at Yulara (0.413 against 0.495 with ACI). Binning the calibration by the GFS
clearness index, a quantity known at the issue time, gives the best interval score at both sites
and raises the coverage on variable days (0.82 against 0.77 at Yulara). Coverage remains below
nominal on overcast and variable days (0.62--0.82 at Yulara, depending on the configuration and
the method), a
limitation shared by marginal conformal methods when the conditioning event is defined by the
outcome.

% ------------------------------------------------------------
\subsection{Sensitivity analyses}
\label{subsec:res_sensitivity}

\paragraph{Season-balanced validation.}
With the season-balanced validation (Section~\ref{subsec:partitions}), the test MAE of AHSE is
38.11 instead of 36.53~kW at Yulara without GFS ($p = 0.015$), 30.86 instead of 31.38~kW with GFS
(n.s.), and 29.16 instead of 27.55~kW at NIST without GFS ($p = 0.05$; 20.33 against 20.39~kW
with GFS, n.s.), all for seed 42. The single tree models improve, because their training weeks
then extend up to the start of the test period, but the selector keeps XGBoost alone and AHSE
loses the benefit of combining members. The chronological protocol is therefore retained.

\paragraph{Foundation model.}
A zero-shot time-series foundation model (Chronos-2, 7-day context) reaches a test MAE of 0.116 in
$k_t$ units at Yulara without GFS, against 0.107 for AHSE, i.e.\ it does not replace a model
trained on the site.

% ------------------------------------------------------------
\subsection{Discussion}
\label{subsec:discussion}

The original claim of this work was that a regime-gated ensemble of six learners outperforms
every single model on Yulara. Three findings revise it. First, the Yulara series used in that
claim is curtailed: forecasting it amounts to forecasting a dispatch rule, and an honest benchmark
must use uncurtailed output or reconstructed available power. Second, once the target is
physical, the ensemble's advantage is real but moderate (7--10\% without weather forecasts at
Yulara, none at NIST), whereas archived numerical weather forecasts, used without look-ahead,
reduce the error by 13\% at Yulara and 26\% at NIST. Third, a statistically guarded selection is
what makes the ensemble safe: it never degraded the best member in our experiments, while the
single-strategy selector of the original version, guarded only by a fixed 0.5\% threshold, was
1.2~kW worse at Yulara without GFS. For an off-grid
mini-grid such as Yulara, the relevant improvement is the GFS-informed forecast together with
calibrated intervals, which size the spinning reserve and the curtailment headroom of the gas
generators directly.

Limitations: one test period per site (10 months at Yulara, 4 months at NIST), a single seed at
NIST, hourly resolution only, and point GFS forecasts rather than an ensemble prediction system.
